\newpage
\section{Lecture-05}
\subsection{Limits of Functions}
\begin{definition}
Let $U \subset \mathbb{R}^n$ be an open subset containing a neighborhood of
$a \in \mathbb{R}^n$, except perhaps for the point $a$ itself.
Suppose $f:U \to \mathbb{R}^m$. We say that
\[
    \lim_{x\to a} f(x)=\ell
\]
($f(x)$ approaches $\ell \in \mathbb{R}^m$ as $x$ approaches $a$) if for every
$\varepsilon>0$ there is $\delta>0$ so that
\[
    \|f(x)-\ell\|<\varepsilon
    \qquad\text{whenever}\qquad
    0<\|x-a\|<\delta.
\]
\end{definition}
\subsection{Relation between the limit of functions and its components}
\begin{proposition}
$$\lim_{x\to a} f(x)=\ell$$ if and only if
$$\lim_{x\to a} f_j(x)=\ell_j$$ for all $j=1,\ldots,m$.
\end{proposition}    
\begin{proof}
    See proposition 3.1, Chapter 2 from \cite{shifrin}.
\end{proof}
\begin{example}
Fix a non-zero vector $b\in\mathbb R^n$. Let $f:\mathbb R^n\rightarrow\mathbb R$ be defined by $f(x)=b\cdot x$. Show that $$\lim_{x\to a} f(x)=b\cdot a$$\\
\textbf{Answer: }Fix a nonzero vector $b \in \mathbb{R}^n$. Let $f : \mathbb{R}^n \to \mathbb{R}$ be defined by
\[
f(x)= b \cdot x.
\]
We claim that
\[
\lim_{x \to a} f(x) = b \cdot a
\]
because
\[
|f(x)-b\cdot a|
= |b\cdot x - b\cdot a|
= |b\cdot (x-a)|
\leq \|b\|\,\|x-a\|,
\]
by the Cauchy-Schwarz Inequality. Thus, given
$\varepsilon > 0$, if we take
\[
\delta = \frac{\varepsilon}{\|b\|},
\]
then whenever $0 < \|x-a\| < \delta$, we have
\[
|f(x)-b\cdot a| < \|b\| \frac{\varepsilon}{\|b\|} = \varepsilon,
\]
as needed.

\end{example}

\begin{example}
Let $f:\mathbb R^n\rightarrow\mathbb R$ be defined by $f(x)=\|x\|^2$. Show that $$\lim_{x\to a} f(x)=\|a\|^2$$\\
\textbf{Answer: }We have
\[
f(x)=\|x\|^2=x\cdot x,
\]
so
\[
f(x)-\|a\|^2=\|x\|^2-\|a\|^2=x\cdot x-a\cdot a.
\]
Now
\[
x\cdot x-a\cdot a=(x-a)\cdot(x+a),
\]
and therefore
\[
|f(x)-\|a\|^2|
= |(x-a)\cdot(x+a)|
\le \|x-a\|\,\|x+a\|
\]
by the Cauchy--Schwarz inequality.

If $\|x-a\|<1$, then
\[
\|x\|\le \|x-a\|+\|a\|<1+\|a\|,
\]
so
\[
\|x+a\|\le \|x\|+\|a\|<1+2\|a\|.
\]
Hence, whenever $\|x-a\|<1$,
\[
|f(x)-\|a\|^2|
\le (1+2\|a\|)\|x-a\|.
\]

Given $\varepsilon>0$, choose
\[
\delta=\min\left\{1,\frac{\varepsilon}{1+2\|a\|}\right\}.
\]
Then if $0<\|x-a\|<\delta$, we have
\[
|f(x)-\|a\|^2|
\le (1+2\|a\|)\|x-a\|
< (1+2\|a\|)\delta
\le \varepsilon.
\]
Thus,
\[
\lim_{x\to a} f(x)=\|a\|^2.
\]
\end{example}

\begin{example}
Let $f:\mathbb R^2\setminus\{0\}\rightarrow\mathbb R$ be defined by $f\left(\begin{bmatrix}
x\\y
\end{bmatrix}\right)=\frac{x^2y}{x^2+y^2}$. Does $\lim_{x\to 0} f(x)$ exist?
\end{example}

\begin{example}
Let $f:\mathbb R^2\setminus\{0\}\rightarrow\mathbb R$ be defined by $f\left(\begin{bmatrix}
    x\\y
\end{bmatrix}\right)=\frac{x^2}{x^2+y^2}$. Does $\lim_{x\to 0} f(x)$ exist?
\end{example} 

\begin{example}
Let $f:\mathbb R^2\setminus\{0\}\rightarrow\mathbb R$ be defined by $f\left(\begin{bmatrix}
    x\\y
\end{bmatrix}\right)=\frac{xy}{x^2+y^2}$. Does $\lim_{x\to 0} f(x)$ exist?
\end{example} 
These examples are taken from the textbook \cite{shifrin}, Examples 3,4 from Chapter 2, section 3.
\begin{proposition}
Suppose $U\subset\mathbb R^n$ is open and $f:U\rightarrow\mathbb R^m$. Then $f$ is continuous at $a\in U$ if and only if for every sequence $\{x_k\}$ of points in $U$ converging to $a$, the sequence $\{f(x_k)\}$ converges to $f(a)$.
\end{proposition}
\begin{proof}
    I have uploaded the recorded video of the proof of this proposition on YouTube (\url{https://www.youtube.com/watch?v=wCV83GlEAGwURL}) which contain the detailed proof of this proposition.

\end{proof}  